\section{Solution Overview}

FalconVQA processes a recording in five stages: the video is partitioned into chunks,
each chunk is analysed, the resulting records are indexed, video-level aggregates are
derived from those records, and the indexed material is then queried by search or by
question. Most stages are parameterised rather than fixed, because surveillance
footage varies too widely for a single configuration to be appropriate in all cases.
This section states what each stage does; Section~5 specifies the models, parameters and
outputs of each in full.

\optfigure{falcon-vqa-arch.png}{FalconVQA system architecture: client and agent tier,
core analysis service, persistence, deployment targets, and the five-phase processing
pipeline.}{1.0}{fig:arch}

\subsection{Partitioning the Recording}

The recording is cut where its content changes. Four boundary detectors contribute
evidence for a cut: a change of speaker, a period of silence, a hard scene cut, and a
shift in the visual content of the frame. Their scores are fused under a weighting, and
the strongest peaks of the fused signal are taken as boundaries.

Where a signal produces no events at all, its weight is redistributed across the
signals that did fire. This behaviour is essential rather than incidental: fixed-camera
footage contains no hard cuts and frequently no speech, and without redistribution the
entire recording is returned as a single chunk. Two further modes are provided for
cases in which fusion is not wanted: a caller-supplied weighting, or fixed spans of
$N$ seconds with no detectors run at all.

The chunking configuration is stored alongside the resulting records, so that a single
recording may be indexed both coarsely and finely at the same time, with a query
resolving against whichever indexing is appropriate.

\subsection{Analysing Each Chunk}

The analyzers to be run are selected at upload time, because their costs differ by
orders of magnitude. Six analyzers are provided: scene description, per-person
description, object detection, on-screen text recognition, transcription, and
speaker-attributed transcription. Transcription and speaker-attributed transcription
are mutually exclusive, as the latter is a strict superset of the former.

The vision--language model is the dominant cost, so inexpensive detectors are run first
and act as gates. If YOLO detects no object in a frame and EasyOCR finds no text, that
frame is never submitted to the vision model. Frames are additionally deduplicated: a
frame is retained only if it differs sufficiently from the last frame retained. The
similarity threshold is adaptive with an absolute ceiling, a configuration arrived at
by measurement: a fixed threshold retained a single frame on a static camera, and a
purely relative threshold began selecting compression noise. Frames are read
sequentially rather than by seeking, which was measured to be approximately
$22\times$ faster.

\subsection{Indexing and Aggregation}

Each chunk is stored as a single point carrying five separately embedded vectors rather
than one: the complete flattened record, the scene description, the people, the
actions, and the objects. A query concerning a person's appearance is therefore matched
against a short description of a person rather than competing against everything else
recorded for that chunk. Embeddings are computed locally, and each point carries a
typed payload against which structured filters are applied.

Twelve aggregators then run over the saved analyzer output, not over the video,
in an order derived from their declared dependencies. They produce the video-level
material: statistics, novelty, speaker statistics, summary, chapters, events, named
entities, sentiment, linked person entities, linked object entities, entity timelines
and co-occurrence. Because they consume stored records, re-running an aggregator incurs
no re-analysis; results are cached and recomputed only on demand or when the analyzer
set changes.

\subsection{Search and Question Answering}

Searches may be restricted by time, objects, tags, people and other typed filters, and
return one of three levels of response detail, selected according to whether the
consumer is a human reader or a token-metered agent.

Questions follow a different path. Rather than searching chunks and synthesising from
whatever is returned, the question is routed to the video-level aggregates capable of
answering it: the entity narratives for a question about a person, novelty for a
question about anomalies, statistics for a question about activity levels.

In both cases the operator supplies a natural-language description and receives ranked
moments carrying \texttt{mm:ss} timecodes. Selecting a result begins playback of the
recording at that point.
